\documentclass[11pt,a4paper]{article}
\usepackage[margin=1in]{geometry}
\usepackage{xcolor}
\usepackage{amsmath}
\usepackage[hidelinks]{hyperref}
\usepackage{enumitem}
\usepackage{titlesec}

\definecolor{revblue}{RGB}{20,60,120}
\newcommand{\comment}[1]{\par\medskip\noindent\textcolor{revblue}{\textbf{Comment.} \textit{#1}}\par\smallskip}
\newcommand{\response}{\par\noindent\textbf{Response.}\ }
\newcommand{\changed}[1]{\par\smallskip\noindent\textbf{Changes:} #1\par}
\newcommand{\TBD}{\textcolor{red}{\textbf{[TBD]}}}
\newcommand{\etal}{\textit{et al.}}

\titleformat{\section}{\large\bfseries}{\thesection.}{0.5em}{}

\title{\vspace{-2em}Response to Reviewers\\[0.3em]
\large Manuscript: \textit{Matryoshka Representation Learning for
Adaptive-Fidelity ECG Classification}\\
\normalsize IEEE Journal of Biomedical and Health Informatics}
\author{}
\date{}

\begin{document}
\maketitle
\vspace{-3em}

\section*{Summary of changes}

We thank both reviewers for reviews that identified real weaknesses rather
than cosmetic ones. Two criticisms in particular changed how we understand our
own result, and we have restructured the paper accordingly rather than
patching around them.

\begin{enumerate}[leftmargin=*,itemsep=3pt]
\item \textbf{The wearable-to-cloud framing was not supported and has been
withdrawn.} The paper now claims adaptive embedding compression on clinical
twelve-lead ECG. The device-tier mapping is presented as arithmetic, not as a
validated deployment result, and the title, abstract, introduction and
conclusion have been rewritten accordingly. We have added reduced-lead and
artefact-corruption experiments that bound the wearable question without
pretending to settle it.

\item \textbf{The computational argument was wrong in kind, not merely in
precision.} Reviewer 1's observation that the backbone dominates inference
cost is correct and is fundamental to the method: prefix truncation occurs
after the encoder has run, so latency is flat in $d$ by construction. We now
measure this on real hardware and state it explicitly as a limitation of MRL.
The previous device-tier latency table, produced by scaling one GPU
measurement by a ratio of nominal device throughputs, has been deleted.

\item \textbf{Fixed-dimension baselines with the same backbone have been
added} at every nesting dimension and every seed. We further note that our
previous headline comparison---an Inception1D MRL model against XResNet1D-101
fixed baselines---confounded nesting with backbone choice and could not
support the conclusion we drew from it.

\item \textbf{Uncertainty quantification throughout}: five seeds per
condition, bootstrap confidence intervals, DeLong tests, Holm correction, and
an equivalence test with a pre-registered margin for the central invariance
claim.

\item \textbf{The clinical feature hierarchy is now measured}, by ridge
probing of eleven physiological descriptors against each embedding prefix. We
commit to reporting the outcome whichever way it falls, including withdrawing
the hierarchy table if the probes do not support it.

\item \textbf{Cross-dataset evaluation} on CPSC-2018, Chapman--Shaoxing and
Georgia, with an explicit statement that the SNOMED-CT to superclass mapping
is lossy.
\end{enumerate}

We also corrected four issues the reviewers did not raise but which we found
while preparing the revision, listed in Section~\ref{sec:selffound}.

\medskip
\noindent\textcolor{red}{\textbf{Note to co-authors:} numbers marked \TBD{} in
the revised manuscript are populated automatically from
\texttt{scripts/aggregate\_results.py} once the experiment matrix has been
run. This letter should be finalised after that point, and the outcome-dependent
passages below (marked $\dagger$) resolved to what was actually measured.}

\section{Reviewer 1}

\subsection*{Major comment 1 --- Wearable claim not supported by PTB-XL}

\comment{The evaluation is performed only on PTB-XL\ldots substantially
different from wearable ECGs, ambulatory recordings, noisy signals, and
single-lead settings\ldots Without one of these additions, the deployment
claim should be recalibrated and framed more narrowly as adaptive embedding
compression on PTB-XL.}

\response We accept this in full, and have taken the recalibration option
seriously rather than treating it as a fallback.

The title no longer contains ``wearable-to-cloud''. Section~I now includes an
explicit subsection, \textit{Scope, and what this paper does not claim},
separating three claims that the previous version conflated: embedding
compressibility (demonstrated), deployment mapping (arithmetic), and on-device
wearable validation (not established). The abstract states the last of these
as an open question.

We have additionally added the strongest proxies available without new data
collection:
\begin{itemize}[leftmargin=*,itemsep=2pt]
\item \textbf{Cross-dataset transfer} to CPSC-2018, Chapman--Shaoxing/Ningbo
and Georgia (Section~V-E1). We report these with the caveat that the
SNOMED-CT to superclass mapping is many-to-one and lossy, so they measure
concept transfer under shift rather than identical-task performance.
\item \textbf{Reduced-lead evaluation} down to single-lead input, in both
zero-shot and natively-trained regimes (Section~V-E2).
\item \textbf{Artefact robustness} across baseline wander, EMG, electrode
motion and mains interference at SNRs from 20\,dB to 0\,dB (Section~V-E3).
\end{itemize}

We are explicit that these are proxies: synthesised artefacts injected into
clinically acquired signals are not dry-electrode wearable recordings, and
masking leads is not a single-lead acquisition geometry. Regarding the
suggestion of HEEDB, MIMIC-IV-ECG or CODE-15\%, our released pipeline supports
these corpora directly; we have not exercised them because they require
credentialed access, and we identify scale as an explicit limitation
(Section~VI-D).

\changed{Title; abstract; new Section~I-A; new Sections~V-E1--V-E3;
Figs. (leads, robustness); Table~(external); Section~VI-B, VI-D.}

\subsection*{Major comment 2a --- Computational benefit not established}

\comment{Reducing the embedding dimensionality reduces representation size,
but the full backbone still has to be executed and appears to dominate
inference cost. Therefore, claims about latency, memory, and energy advantages
require real-device benchmarking or a more careful limitation of the claims.}

\response This is the most important comment in either review, and the
reviewer is right about the mechanism, not merely about the evidence. Prefix
truncation is a pointer offset applied \emph{after} the encoder has run;
latency, energy and peak activation memory are therefore invariant to $d$ by
construction. No amount of additional measurement would change this.

We have done both things the reviewer offers as alternatives. We now measure
on real hardware---GPU with CUDA-event timing and NVML energy integration,
multi- and single-threaded CPU, dynamically quantised CPU, and ONNX
Runtime---and we have rewritten the efficiency claim. Section~V-C now states
that nesting reduces (i) embedding storage and transmission, by up to
32$\times$, and (ii) the number of maintained model artefacts, with the
associated validation and certification burden; and that it does \emph{not}
reduce inference latency, energy or activation memory. Table (related
work) now carries a ``compute saving'' column in which our own method is
marked \textit{No}.

The previous Table~III has been deleted. Its Cortex-M4 column was an
analytical extrapolation from an Apple MPS measurement scaled by nominal
device throughput, which we should not have presented in a table of measured
quantities.

\changed{Section~V-C rewritten; new measured Table~(compute) and
Fig.~(latency); old Table~III removed; Section~VI-B; related-work table
column added; \texttt{scripts/benchmark\_hardware.py} released.}

\subsection*{Major comment 2b --- Fixed-dimension Inception1D baselines}

\comment{Since the best proposed model uses Inception1D, fixed-dimension
Inception1D baselines should be included.}

\response Added, at all six dimensions and all five seeds, for both backbones.

We want to acknowledge the deeper issue this exposes. Our previous headline
claim---that the MRL model ``outperforms six independently trained
34M-parameter baselines''---compared an Inception1D MRL model against
XResNet1D-101 fixed baselines. Since Inception1D also outperformed
XResNet1D-101 \emph{under MRL}, that comparison confounded nesting with
backbone choice and could not establish what we claimed. The matched baselines
now separate the two factors, and we have removed the inference that the
multi-granularity loss acts as a regulariser, which rested on the same
confounded comparison (Section~VI-C).

\changed{Table~(main) restructured; Section~V-B; Section~VI-C;
\texttt{run\_all.sh} Stage~3.}

\subsection*{Major comment 2c --- Repeated runs and statistical testing}

\comment{Repeated runs, confidence intervals, or statistical testing are also
needed, because some reported differences are small.}

\response Every condition is now trained with five seeds and reported as mean
$\pm$ SD. Test-set uncertainty uses bootstrap over recordings ($B$=2{,}000);
model comparisons use paired bootstrap and DeLong's test with Holm correction
across the six dimensions.

One methodological point we would highlight: our central claim is that
performance is \emph{stable} across dimensions, which is a claim of no
meaningful difference. A non-significant superiority test cannot establish
that. We therefore pre-register an equivalence margin of $\pm$0.01 macro AUC
(roughly half the seed-to-seed SD of a PTB-XL classifier) and require the
whole confidence interval of the difference to fall inside it
(Section~IV-E, Table~(stats)).

\changed{New Section~IV-E; Table~(stats); all tables report SD; seed bands in
all figures; \texttt{mecg/analysis/stats.py}.}

\subsection*{Minor comment 1 --- Terminology}

\comment{The terminology around ``wearable'', ``IoT'', and ``deployment''
should be made more precise, distinguishing between conceptual deployment
mapping, actual on-device inference, and validation on wearable data.}

\response We have adopted exactly this three-way distinction and made it
structural: it is enumerated in Section~I-A, maintained consistently in
Sections~III-F, V-C and V-E, and revisited in the discussion. ``IoT'' has
largely been removed, since it was doing rhetorical rather than technical
work. We now write ``deployment mapping'' for the arithmetic,
``on-device inference'' for execution on target hardware, and ``wearable
validation'' for evaluation on wearable-acquired data, and we state which of
the three each claim belongs to.

\section{Reviewer 2}

\subsection*{Comment 1 --- Does not outperform stronger PTB-XL classifiers}

\comment{The proposed model does not outperform stronger reported PTB-XL
classifiers and its contribution is mainly adaptive compression rather than
improved ECG diagnosis.}

\response We accept the framing point and have acted on it, but the factual
premise has changed during revision. Having corrected the normalisation defect
described in Section~\ref{sec:selffound}, our Inception1D-MRL model now
attains 0.9269 macro AUC---above the 0.921 reported for the same backbone by
Strodthoff \etal. We do not build the paper's contribution on that, and
Section~V-G states plainly that adaptivity rather than accuracy is the
contribution. But we can now demonstrate the flat accuracy--dimension profile
at benchmark-level performance, which is a stronger claim than the original
submission supported. Section~V-G presents the comparison with an
explicit statement that we make no accuracy claim, and that published
non-adaptive systems each serve one operating point. Section~I states that we
decline the superiority framing, and the abstract characterises the
contribution as concerning the \emph{shape} of the accuracy--dimension curve
rather than its height.

We also note---since it bears on the reviewer's point---that our absolute
performance is below the Strodthoff \etal{} benchmark for the same
architectures. We identified a probable cause during revision
(Section~\ref{sec:selffound}, item 1) and have corrected it; whether the gap
closes will be evident in the revised results.

\subsection*{Comment 2 --- Missing fixed Inception1D baseline}

\response Addressed above (Reviewer 1, major 2b).

\subsection*{Comment 3 --- Deployment claims not validated on hardware}

\response Addressed above (Reviewer 1, major 2a). Measurements are now real;
the claim has been narrowed to what measurement supports.

\subsection*{Comment 4 --- Clinical hierarchy is speculative}

\comment{The proposed clinical hierarchy of embedding dimensions is
speculative and not validated by probing or interpretability analysis.}

\response Accepted. The previous Table~II asserted a mapping from nesting
dimensions to clinical features on the basis of a gradient-flow argument that
establishes only that leading coordinates receive more gradient signal---not
that they encode any \emph{particular} feature. The inference did not follow.

We now test it. Eleven physiological descriptors (RR statistics, QRS duration,
QT, ST level, T-wave amplitude, R and P amplitude, frontal axis) are extracted
from the raw waveforms; ridge probes are fitted on training-split embeddings
and evaluated out-of-sample; and we record the saturation dimension for each
descriptor. Section~V-D reports the rank correlation between hypothesised and
measured ordering, plus prefix CKA, effective rank and the variance profile.

$\dagger$ We state in the manuscript, before seeing the outcome, what each
result would mean: if saturation order tracks the hypothesis, the nesting
structure has a physiological interpretation; if all descriptors saturate at
the smallest prefix, the parsimonious conclusion is that the diagnostic
manifold is simply low-dimensional, which explains the flat AUC profile
equally well but licenses no claim about what particular coordinate ranges
encode. In that case the hierarchy table will be withdrawn rather than
reinterpreted.

\changed{Old Table~II removed pending the probing result; new Section~V-D;
Table~(probe); Fig.~(probing); \texttt{mecg/analysis/probing.py},
\texttt{features.py}.}

\subsection*{Comment 5 --- Limited to PTB-XL five-class}

\response Cross-dataset evaluation added (Reviewer 1, major 1). On finer
granularity: we have not added the 44-class task, and we explain why this is
more than a resource question. If the flat profile is explained by low
intrinsic dimensionality, the sufficient prefix should \emph{grow} with task
granularity---so the 44-class result is a genuine test of the mechanism rather
than a routine extension, and we would rather flag it as the sharpest open
prediction (Section~VI-D, VI-E) than report it hastily. Our preprocessing
pipeline supports it via \texttt{--task subdiagnostic}.

\subsection*{Comment 6 --- No uncertainty analysis}

\response Addressed above (Reviewer 1, major 2c).

\section{Issues we identified ourselves}
\label{sec:selffound}

\begin{enumerate}[leftmargin=*,itemsep=4pt]
\item \textbf{Normalisation destroyed voltage information.} The original
pipeline applied per-record, per-lead $z$-scoring, dividing each lead by its
own standard deviation. This removes \emph{absolute} amplitude---which is the
diagnostic criterion for ventricular hypertrophy (Sokolow--Lyon, Cornell).
HYP was by far our worst class (AUC 0.811 against 0.93 for NORM), and we
believe this is why. Normalisation now uses training-split statistics; the
comparison is retained as an ablation. This may also account for part of the
gap to the published PTB-XL benchmark noted by Reviewer 2.

\item \textbf{Weight decay applied to normalisation parameters.} The
optimiser grouped parameters by the substring test \texttt{'bn' in name},
which missed every BatchNorm nested inside a convolutional block. Now grouped
by \texttt{param.ndim <= 1}.

\item \textbf{F1 evaluated at a fixed 0.5 threshold.} For an imbalanced
multi-label task this substantially understates F1 and explains most of the
apparent inconsistency between our AUC ($\approx$0.90) and F1 ($\approx$0.62).
Thresholds are now tuned per class on validation and frozen for test.

\item \textbf{Paper--code mismatch and citation errors.} The manuscript
described the Inception1D backbone as using squeeze-and-excitation attention;
the released implementation did not. SE is now genuinely enabled. We also
corrected citations: XResNet1D was attributed to He \etal{} (ResNet) and
Inception1D to Szegedy \etal{} (GoogLeNet), where the relevant references are
Strodthoff \etal{} and Ismail Fawaz \etal{} (InceptionTime) respectively.

\item \textbf{SVD baseline leakage.} The basis and classifier were fitted on
the test partition and evaluated on it. The manuscript acknowledged this as
``optimistically biased'' but nonetheless plotted the curve as an upper bound
against honestly evaluated models. Both are now fitted on the training split.

\item \textbf{Model selection on the largest head.} Early stopping used
validation AUC at $d$=512, biasing the checkpoint toward one operating point.
Selection now uses the mean across granularities.

\item \textbf{Figure axis scaling.} The original Pareto figure used a
0.02-AUC vertical window, magnifying sub-0.005 differences into visually
dramatic gaps and suggesting structure that is within seed noise. Axes are no
longer zoomed and seed variability is shown.
\end{enumerate}

\section*{Closing}

Two of these criticisms---the compute mechanism and the confounded
baseline---changed our understanding of what we had shown, and the revision is
correspondingly a restructuring rather than an addition of experiments. We
believe the narrower claim is the more useful contribution, and we are
grateful to both reviewers for pressing on exactly the points that needed it.

\end{document}
